% Generated draft: requires mathematical and source review.
\documentclass[11pt]{article}
\usepackage[margin=25mm]{geometry}
\usepackage{fontspec}
\setmainfont{Latin Modern Roman}
\usepackage{amsmath,amssymb,mathtools,unicode-math}
\setmathfont{Latin Modern Math}
\usepackage{enumitem,needspace,xurl,hyperref}
\setlength{\parindent}{0pt}
\setlength{\parskip}{5pt}
\setlength{\emergencystretch}{3em}
\newcommand{\N}{\mathbb{N}}
\newcommand{\Z}{\mathbb{Z}}
\newcommand{\Q}{\mathbb{Q}}
\newcommand{\R}{\mathbb{R}}
\newcommand{\C}{\mathbb{C}}
\newcommand{\F}{\mathbb{F}}
\newcommand{\A}{\mathbb{A}}
\newcommand{\PP}{\mathbb{P}}
\newcommand{\E}{\mathbb{E}}
\newcommand{\eps}{\varepsilon}
\newcommand{\dd}{\,\mathrm{d}}
\begin{document}
\section*{UnsolvedMath ID: 30000256}
\textbf{Source problem code:} catalogue entry
\subsection{Definitions and mathematical statement}
Let \(\Omega\subset\R^N\) be the bounded smooth domain intended in the source, and let \(\nu\) be its outward unit normal. Consider a family of positive solutions, as \(\varepsilon\downarrow0\), of
\[
 \varepsilon^2\Delta u_\varepsilon-u_\varepsilon+u_\varepsilon^p=0
 \quad\text{in }\Omega,\qquad
 \partial_\nu u_\varepsilon=0\quad\text{on }\partial\Omega.
\]
The energy assumption is that there exist \(C>0\) and \(\varepsilon_0>0\), independent of \(\varepsilon\), such that
\[
 E_\varepsilon[u_\varepsilon]\leq C\varepsilon^{N-m}
 \qquad(0<\varepsilon<\varepsilon_0).
\]
The historical question is whether this forces the energy-concentration set of the family to have dimension exactly \(m\). The exponent \(p\), the permitted values of \(N,m\), the solution regularity and the functional \(E_\varepsilon\) are not defined in the supplied record, so no subcriticality or other missing hypothesis is inserted.

For example, a conventional variational functional for the displayed equation, when the terms are defined and \(p\ne-1\), is
\[
 \mathcal E_\varepsilon(u)=\int_\Omega
 \left(\frac{\varepsilon^2}{2}|\nabla u|^2+\frac12u^2-
 \frac{1}{p+1}u^{p+1}\right)\,dx.
\]
This is only a conditional notation reference: the record does not identify \(E_\varepsilon\) with \(\mathcal E_\varepsilon\) or specify any rescaling. That identification matters to the stated energy scale. An energy-concentration set likewise needs a definition, such as one derived from a specified normalized energy measure along a specified subsequence; no particular measure, support or pointwise blow-up criterion is given. Its location may include the boundary under the Neumann condition. The meaning of dimension, including Hausdorff dimension versus dimension of a smooth concentration locus, is also absent. An upper energy bound is the stated assumption; a matching lower bound or regularity of the concentration set is not added to make an exact-dimension conclusion plausible.

There is an explicit status conflict in the supplied record. The verification field says the official report poses this precise question and then answers it negatively by constructing many-spike solutions, retaining it as a verified historical problem. In contrast, the 2026-08-21 literature assessment calls the general claim unresolved and the registry status is open. The draft therefore records the question as historical and attributes the reported negative answer to that verification note; it does not promote the affirmative implication to a currently open conjecture. The actual many-spike construction, its parameter hypotheses and its concentration-set calculation are not supplied and are not invented here.
\subsection{Short English statement}
For positive Neumann solutions of the stated singularly perturbed equation, does an upper energy bound of order \(\varepsilon^{N-m}\) force an \(m\)-dimensional concentration set? The source-verification note reports a negative answer by many-spike solutions, contradicting the older open-status assessment. The historical question is retained with that conflict and the missing energy and concentration definitions explicit.
\subsection{Sources}
\begin{itemize}
\item Local UnsolvedMath registry, permanent ID 30000256, source code catalogue entry, ``Dimensions of Energy-Concentration Sets,'' including original question, expanded PDE and boundary conditions, and the verification note reporting a negative many-spike answer.
\item \emph{Topological and Variational Methods for Differential Equations (2006)}, Oberwolfach Reports, pp.\ 1601--1678, DOI 10.4171/owr/2005/29: \url{https://doi.org/10.4171/owr/2005/29}. The citation uses 2006, while the DOI and proposed-year metadata use 2005.
\item Supplied reference ``Druet, Lin--Ni problem for mean-convex domains'': \url{https://arxiv.org/abs/1103.3811}.
\item Literature triage checked 2026-08-21; record updated 2026-09-16. Its unresolved assessment conflicts with the source-verification note's negative-answer report. No external source has been visited or independently checked here.
\end{itemize}
\end{document}
